\chapter{Discussion}
\label{ch:discussion}

Our project investigated an interpretable sparse dictionary learning framework for abnormality classification and localization in 3D chest CT scans. A key contribution is the frozen dictionary approach, where normal anatomical patterns are learned first and abnormalities are learned separately. This helps isolate pathological features and improves interpretability.

Unlike many stateoftheart deep learning models that rely on posthoc explainability methods such as GradCAM, the proposed LC-KSVD framework provides inherent explainability. The generated predictions can be directly linked to dictionary atoms and reconstructed voxelspace regions through atom backprojection, allowing more transparent clinical interpretation.

The experimental observations show that the framework can successfully classify and localize pulmonary nodules, Ground Glass Opacity (GGO), and consolidation patterns from 3D chest CT scans while maintaining competitive performance. The proposed objectives of implementing an interpretable dictionary learning pipeline, generating localization maps, and evaluating the framework against existing approaches were successfully achieved.

In realworld clinical applications, this approach could support radiologists by providing both diagnostic predictions and anatomically grounded explanations, improving trust and transparency in AIassisted diagnosis systems.

However, the proposed method has several limitations. The framework is currently limited to three abnormalities and does not yet provide a fully unified endtoend classification and segmentation architecture. In addition, although the explainability is stronger, the classification accuracy may still be lower than largescale deep learning foundation models. Future work will focus on improving performance, expanding abnormality coverage, and integrating more advanced hybrid architectures.
